\subsection{\texorpdfstring{Relaciones entre las propiedades de un anillo local noetheriano $A$ y
las de un anillo cociente $A/tA$.}{Relaciones entre las propiedades de un anillo local noetheriano A y las de un anillo cociente A/tA}}

Ya se han visto en (3.4) relaciones entre las propiedades de $A$ y las de $A/tA$
relativas a los ideales primos asociados, así como a las propiedades de ser
íntegro o reducido. En esta sección se dan otras relaciones entre las
propiedades de estos anillos, ligadas a las nociones de dimensión y profundidad.

\begin{proposition}[5.12.1]
\label{IV.2.s5.12.1-fr}
Sean $A$ un anillo local noetheriano, $X=\operatorname{Spec}(A)$, $t$ un
elemento de $A$ que forma parte de un sistema de parámetros \textup{(\textbf 0, 16.3.6)}, $X_0$
el subespacio cerrado $V(t)$ de $X$, y $X_1$ el abierto complementario
$X\setmin X_0$. Sea $Z$ una parte de $X$ tal que toda especialización de un punto de
$Z$ pertenezca a $Z$. Se supone que $X$ es biequidimensional \textup{(lo que ocurrirá
si $A$ es equidimensional y cociente de un anillo local regular (\textbf{0}, 16.5.12))}.
Entonces, si se pone $Z_0=Z\cap X_0$, $Z_1=Z\cap X_1$, se tiene
\begin{equation}
\operatorname{codim}(Z_0,X_0)\leq\operatorname{codim}(Z,X)
\leq\operatorname{codim}(Z_1,X_1).
\tag{5.12.1.1}\label{IV.2.s5.12.1.1-fr}
\end{equation}
\end{proposition}

La segunda desigualdad resulta de la definición (5.1.3); para demostrar la
primera, es posible limitarse a demostrarla cuando $Z$ es \emph{cerrado}: en efecto, por
hipótesis, $Z$ es unión de partes cerradas $Z_\alpha$, y si se tiene
$\operatorname{codim}(X_0\cap Z_\alpha,X_0)
\leq\operatorname{codim}(Z_\alpha,X)$ para todo $\alpha$, se tendrá también
$\operatorname{codim}(Z_0,X_0)
=\inf_\alpha(\operatorname{codim}(X_0\cap Z_\alpha,X_0))
\leq\inf_\alpha(\operatorname{codim}(Z_\alpha,X))
=\operatorname{codim}(Z,X)$.
Supongamos, pues, $Z$ cerrado, de modo que se tiene
\[
\operatorname{codim}(Z,X)=\dim(X)\setmin\dim(Z)
\]

\oldpage[IV-2]{127}

en virtud de (\textbf 0, 14.3.5.1). Por otra parte, $X_0$ es evidentemente catenario,
equidimensional y de dimensión $\dim(X)\setmin 1$ en virtud de (\textbf 0, 16.3.4); se tiene
por tanto también
\[
\operatorname{codim}(Z_0,X_0)=\dim(X_0)\setmin\dim(Z_0)
=\dim(X)\setmin 1\setmin\dim(Z_0)
\]
Pero se tiene $\dim(Z_0)\geq\dim(Z)\setmin 1$ (\textbf 0, 16.3.4), lo que termina de probar la
primera desigualdad (5.12.1.1).

\begin{proposition}[5.12.2]
\label{IV.2.s5.12.2-fr}
Sean $A$ un anillo local noetheriano, $M$ un $A$-módulo de tipo finito, $t$ un
elemento $M$-regular perteneciente al ideal maximal $\mathfrak m$ de $A$, y
$k$ un entero $\geq1$. Se supone que $A$ es un anillo catenario. Entonces, si
$M/tM$ es equidimensional y verifica $(S_k)$, lo mismo ocurre con $M$.
\end{proposition}

Teniendo en cuenta $(\mathbf{Err}_{\mathrm{III}},30)$ y (5.7.3, (vi)), es posible
limitarse al caso en que $\operatorname{Supp}(M)=\operatorname{Spec}(A)=X$.
Pongamos $X_0=V(t)=\operatorname{Spec}(A/tA)$; como $M/tM$ verifica $(S_1)$ por
hipótesis, no tiene ideales primos asociados inmersos (5.7.5). Aplicando (3.4.4)
al punto cerrado de $X$ (y de todo subpreesquema cerrado de $X$), se ve que las
componentes irreducibles de $\operatorname{Supp}(M/tM)$ son exactamente las de
$Y_i\cap X_0$, donde $Y_i$ ($1\leq i\leq r$) son las componentes irreducibles
de $X$. Como $t$ es, por hipótesis, $M$-regular, $V(t)$ no contiene ningún
punto maximal de $X$ (3.1.8); cada componente irreducible de $Y_i\cap X_0$ es,
por tanto, de codimensión 1 en $Y_i$ (5.1.8); como $X$ es catenario y todas las
componentes irreducibles de $\operatorname{Supp}(M/tM)$ tienen, por hipótesis,
la misma dimensión, se concluye que todos los $Y_i$ tienen la misma dimensión;
dicho de otro modo, $X$ es equidimensional, y por tanto \emph{biequidimensional},
puesto que $A$ es local y catenario. Para ver que $M$ verifica $(S_k)$,
apliquemos el criterio (5.7.4): hay que comprobar (con las notaciones de
(5.7.4)) que, para todo entero $n\geq0$, se tiene
$\operatorname{codim}(Z_n,X)>n+k$. Ahora bien, la hipótesis sobre $M/tM$ y el
criterio (5.7.4) muestran que se tiene
$\operatorname{codim}(Z_n\cap X_0,X_0)>n+k$ para todo entero $n\geq0$. Pero
toda especialización de un punto de $Z_n$ pertenece también a $Z_n$, como
veremos en (6.11.5)
\footnote{El lector comprobará que (5.12.2) no se utiliza en la demostración
de (6.11.5).}; como $X$ es biequidimensional y $t$ pertenece a un sistema de
parámetros para $M$ (\textbf 0, 16.4.1), y por tanto también para $A$ (pues el anulador
de $M$ es nilpotente en virtud de la hipótesis
$\operatorname{Supp}(M)=\operatorname{Spec}(A)$), la conclusión resulta de
(5.12.1).

\begin{remark}[5.12.3]
\label{IV.2.s5.12.3-fr}
Si, en el enunciado de (5.12.2), se supone solamente que $M/tM$ es
equidimensional, no se sigue necesariamente que $M$ sea equidimensional. Por
ejemplo, sean $k$ un cuerpo, $B$ el anillo de polinomios $k[T,U,V]$ en tres
indeterminadas, $C$ el anillo local de $B$ en el ideal maximal $BT+BU+BV$ de
$B$, y sean $\mathfrak p=CU$, $\mathfrak q=CV+C(T+U)$; consideremos entonces
el anillo local $A=C/\mathfrak p\mathfrak q$ (geométricamente, si $X$ es el
subpreesquema cerrado del espacio afín de tres dimensiones sobre $k$, formado
por un plano y una recta que corta dicho plano en un punto $x$, $A$ es el
anillo local de $X$ en el punto $x$). Sean $t,u,v$ las imágenes canónicas de
$T,U,V$ en $A$; es inmediato que $t$ no es divisor de cero en $A$, y $A/tA$
es isomorfo a $C_0/\mathfrak p_0\mathfrak q_0$, donde $C_0$ es el anillo local
de $B_0=k[U,V]$ en el ideal maximal $B_0U+B_0V$, y
$\mathfrak p_0=C_0U$, $\mathfrak q_0=C_0U+C_0V$ (siendo $\mathfrak q_0$, por
tanto, el ideal maximal de $C_0$). Es inmediato que $\operatorname{Spec}(A/tA)$

\oldpage[IV-2]{128}

es irreducible y de dimensión 1, $A/tA$ no verifica $(S_1)$ y
$\operatorname{Spec}(A)$ no es equidimensional.
\end{remark}

\begin{corollary}[5.12.4]
\label{IV.2.s5.12.4-fr}
Sean $A$ un anillo local noetheriano catenario, $t$ un elemento regular
perteneciente al ideal maximal de $A$, y $k$ un entero $\geq1$. Si $A/tA$
verifica $(S_k)$, lo mismo ocurre con $A$.
\end{corollary}

Si $k=1$, el corolario resulta de (3.4.4), dada la interpretación de la
condición $(S_1)$ (5.7.5). Si $k\geq2$, del criterio de Hartshorne (5.10.9) se
deduce que $A/tA$ es equidimensional, y se puede aplicar (5.12.2).
